% =============================================================================
% Chapter 7: Transformer Blocks
% =============================================================================
\chapter{Transformer Blocks}
\label{ch:transformer-blocks}

\begin{objectives}
    \item Understand the feed-forward network and SwiGLU activation
    \item Learn about RMSNorm and why it's preferred over LayerNorm
    \item Master the pre-normalization pattern
    \item Understand residual connections and gradient flow
    \item See how components combine into a complete transformer block
\end{objectives}

% -----------------------------------------------------------------------------
\section{Anatomy of a Transformer Block}
% -----------------------------------------------------------------------------

A transformer block combines self-attention with a feed-forward network, connected by residual connections and normalization layers.

\begin{figure}[H]
\centering
\begin{tikzpicture}[
    box/.style={rectangle, draw, rounded corners, minimum width=3cm, minimum height=0.8cm, align=center},
    norm/.style={box, fill=yellow!20},
    attn/.style={box, fill=orange!20},
    ffn/.style={box, fill=green!20},
    add/.style={circle, draw, minimum size=0.6cm, fill=gray!10},
    arrow/.style={-{Stealth[length=2mm]}, thick},
]
    % Input
    \node[box, fill=blue!10] (input) at (0, 0) {Input $x$};

    % First norm + attention
    \node[norm] (norm1) at (0, -1.2) {RMSNorm};
    \node[attn] (attn) at (0, -2.4) {Multi-Head\\Self-Attention};

    % First residual add
    \node[add] (add1) at (0, -3.6) {$+$};

    % Second norm + FFN
    \node[norm] (norm2) at (0, -4.8) {RMSNorm};
    \node[ffn] (ffn) at (0, -6) {Feed-Forward\\(SwiGLU)};

    % Second residual add
    \node[add] (add2) at (0, -7.2) {$+$};

    % Output
    \node[box, fill=blue!10] (output) at (0, -8.4) {Output};

    % Main flow
    \draw[arrow] (input) -- (norm1);
    \draw[arrow] (norm1) -- (attn);
    \draw[arrow] (attn) -- (add1);
    \draw[arrow] (add1) -- (norm2);
    \draw[arrow] (norm2) -- (ffn);
    \draw[arrow] (ffn) -- (add2);
    \draw[arrow] (add2) -- (output);

    % Residual connections
    \draw[arrow, dashed, gray] (input.west) -- ++(-0.8, 0) -- ++(0, -3.6) -- (add1.west);
    \draw[arrow, dashed, gray] (add1.west) -- ++(-0.8, 0) -- ++(0, -3.6) -- (add2.west);

    % Labels
    \node[right=0.5cm of add1, font=\footnotesize, gray] {Residual};
    \node[right=0.5cm of add2, font=\footnotesize, gray] {Residual};
\end{tikzpicture}
\caption{A complete transformer block with pre-normalization}
\label{fig:transformer-block}
\end{figure}

The block can be expressed as:
\begin{align}
x' &= x + \text{Attention}(\text{RMSNorm}(x)) \\
\text{output} &= x' + \text{FFN}(\text{RMSNorm}(x'))
\end{align}

% -----------------------------------------------------------------------------
\section{RMSNorm: Efficient Normalization}
% -----------------------------------------------------------------------------

\TryLM uses \textbf{RMSNorm} (Root Mean Square Normalization) instead of the original Transformer's LayerNorm.

\subsection{LayerNorm vs RMSNorm}

\textbf{LayerNorm} computes mean and variance:
\begin{equation}
\text{LayerNorm}(x) = \frac{x - \mu}{\sqrt{\sigma^2 + \epsilon}} \cdot \gamma + \beta
\end{equation}

\textbf{RMSNorm} only computes the RMS (root mean square):
\begin{equation}
\text{RMSNorm}(x) = \frac{x}{\sqrt{\frac{1}{d}\sum_{i=1}^{d} x_i^2 + \epsilon}} \cdot \gamma
\end{equation}

\subsection{Why RMSNorm?}

\begin{enumerate}
    \item \textbf{Simpler}: No mean computation, no bias parameter
    \item \textbf{Faster}: Fewer operations per normalization
    \item \textbf{Equally effective}: Empirically matches LayerNorm performance
\end{enumerate}

\subsection{Implementation}

\begin{lstlisting}[style=python, caption={RMSNorm from model.py}]
class RMSNorm(nn.Module):
    def __init__(self, dim: int, eps: float = 1e-6) -> None:
        super().__init__()
        self.eps = eps
        self.weight = nn.Parameter(torch.ones(dim))  # Learned scale

    def forward(self, x: Tensor) -> Tensor:
        # Compute RMS: sqrt(mean(x^2))
        rms = torch.sqrt(x.pow(2).mean(-1, keepdim=True) + self.eps)

        # Normalize and scale
        return x / rms * self.weight
\end{lstlisting}

\begin{keypoint}
RMSNorm has only \textbf{one learnable parameter} per dimension (the scale $\gamma$), while LayerNorm has two (scale and shift). For $d_{\text{model}} = 384$, each RMSNorm layer has 384 parameters.
\end{keypoint}

% -----------------------------------------------------------------------------
\section{Feed-Forward Network with SwiGLU}
% -----------------------------------------------------------------------------

The feed-forward network (FFN) is the ``processing'' layer that follows attention. \TryLM uses \textbf{SwiGLU}, a gated activation function.

\subsection{Standard FFN}

The original Transformer FFN:
\begin{equation}
\text{FFN}(x) = \text{GELU}(xW_1 + b_1)W_2 + b_2
\end{equation}

\subsection{SwiGLU: Gated Linear Unit with Swish}

SwiGLU uses a gating mechanism:
\begin{equation}
\text{SwiGLU}(x) = \text{Swish}(xW_{\text{gate}}) \odot (xW_{\text{up}})
\end{equation}

where $\odot$ is element-wise multiplication and Swish (also called SiLU) is:
\begin{equation}
\text{Swish}(x) = x \cdot \sigma(x) = \frac{x}{1 + e^{-x}}
\end{equation}

\begin{figure}[H]
\centering
\begin{tikzpicture}[
    box/.style={rectangle, draw, rounded corners, minimum width=1.8cm, minimum height=0.6cm, fill=blue!10, align=center},
    proj/.style={rectangle, draw, minimum width=1.2cm, minimum height=0.6cm},
    op/.style={circle, draw, minimum size=0.5cm, fill=gray!10},
    arrow/.style={-{Stealth[length=2mm]}, thick},
]
    % Input
    \node[box, minimum width=2.5cm] (input) at (0, 0) {Input $x$\\\shape{B, S, 384}};

    % Gate and Up projections
    \node[proj, fill=red!20] (gate) at (-1.5, -1.5) {$W_{\text{gate}}$};
    \node[proj, fill=green!20] (up) at (1.5, -1.5) {$W_{\text{up}}$};

    % Swish activation
    \node[box, fill=orange!20] (swish) at (-1.5, -3) {Swish};

    % Element-wise multiply
    \node[op] (mul) at (0, -4.5) {$\odot$};

    % Down projection
    \node[proj, fill=purple!20] (down) at (0, -6) {$W_{\text{down}}$};

    % Output
    \node[box, minimum width=2.5cm] (output) at (0, -7.5) {Output\\\shape{B, S, 384}};

    % Arrows
    \draw[arrow] (input) -- (-1.5, -0.5) -- (gate);
    \draw[arrow] (input) -- (1.5, -0.5) -- (up);
    \draw[arrow] (gate) -- (swish);
    \draw[arrow] (swish) -- (mul);
    \draw[arrow] (up) -- (1.5, -4.5) -- (mul);
    \draw[arrow] (mul) -- (down);
    \draw[arrow] (down) -- (output);

    % Dimension annotations
    \node[font=\footnotesize, gray, right=0.2cm of gate] {\shape{384, 1024}};
    \node[font=\footnotesize, gray, right=0.2cm of up] {\shape{384, 1024}};
    \node[font=\footnotesize, gray, right=0.2cm of down] {\shape{1024, 384}};
\end{tikzpicture}
\caption{SwiGLU feed-forward network architecture}
\label{fig:swiglu}
\end{figure}

\subsection{Why SwiGLU?}

The gating mechanism allows the network to \textbf{selectively pass information}:
\begin{itemize}
    \item The ``gate'' path learns \textit{what} to let through
    \item The ``up'' path provides \textit{what} to pass
    \item Their product gives controlled information flow
\end{itemize}

Empirically, SwiGLU outperforms GELU/ReLU on language modeling benchmarks.

\subsection{Hidden Dimension}

For SwiGLU, we use $\frac{2}{3}$ of the standard hidden dimension to maintain similar parameter count:

\begin{lstlisting}[style=python, caption={FFN with SwiGLU from model.py}]
class FeedForward(nn.Module):
    def __init__(self, config: ModelConfig):
        super().__init__()

        # SwiGLU uses 2/3 of the standard hidden dim
        hidden_dim = int(2 * config.d_ff / 3)

        # Round to nearest multiple of 256 for efficiency
        hidden_dim = 256 * ((hidden_dim + 255) // 256)

        self.gate_proj = nn.Linear(config.d_model, hidden_dim, bias=False)
        self.up_proj = nn.Linear(config.d_model, hidden_dim, bias=False)
        self.down_proj = nn.Linear(hidden_dim, config.d_model, bias=False)
        self.dropout = nn.Dropout(config.dropout)

    def forward(self, x: Tensor) -> Tensor:
        # SwiGLU: swish(gate) * up
        gate = F.silu(self.gate_proj(x))  # Swish = SiLU
        up = self.up_proj(x)
        hidden = gate * up
        hidden = self.dropout(hidden)
        return self.down_proj(hidden)
\end{lstlisting}

\begin{keypoint}
For \TryLM with $d_{\text{model}} = 384$ and $d_{\text{ff}} = 1536$:
\begin{align*}
\text{hidden\_dim} &= \frac{2}{3} \times 1536 = 1024
\end{align*}
After rounding to 256: hidden\_dim = 1024 (unchanged in this case).
\end{keypoint}

% -----------------------------------------------------------------------------
\section{Pre-Normalization}
% -----------------------------------------------------------------------------

\TryLM uses \textbf{pre-normalization} (Pre-LN): normalize \textit{before} the sublayer, not after.

\subsection{Post-Norm vs Pre-Norm}

\textbf{Post-normalization} (original Transformer):
\begin{equation}
x' = \text{Norm}(x + \text{Sublayer}(x))
\end{equation}

\textbf{Pre-normalization} (\TryLM):
\begin{equation}
x' = x + \text{Sublayer}(\text{Norm}(x))
\end{equation}

\begin{figure}[H]
\centering
\begin{tikzpicture}[
    box/.style={rectangle, draw, rounded corners, minimum width=2cm, minimum height=0.6cm},
    arrow/.style={-{Stealth[length=2mm]}, thick},
]
    % Post-norm
    \node at (-3, 2) {\textbf{Post-Norm:}};
    \node[box, fill=blue!10] (x1) at (-3, 1) {$x$};
    \node[box, fill=orange!20] (sub1) at (-3, -0.3) {Sublayer};
    \node[circle, draw, fill=gray!10] (add1) at (-3, -1.6) {$+$};
    \node[box, fill=yellow!20] (norm1) at (-3, -2.9) {Norm};

    \draw[arrow] (x1) -- (sub1);
    \draw[arrow] (sub1) -- (add1);
    \draw[arrow] (add1) -- (norm1);
    \draw[arrow, dashed, gray] (x1.west) -- ++(-0.5, 0) |- (add1.west);

    % Pre-norm
    \node at (3, 2) {\textbf{Pre-Norm:}};
    \node[box, fill=blue!10] (x2) at (3, 1) {$x$};
    \node[box, fill=yellow!20] (norm2) at (3, -0.3) {Norm};
    \node[box, fill=orange!20] (sub2) at (3, -1.6) {Sublayer};
    \node[circle, draw, fill=gray!10] (add2) at (3, -2.9) {$+$};

    \draw[arrow] (x2) -- (norm2);
    \draw[arrow] (norm2) -- (sub2);
    \draw[arrow] (sub2) -- (add2);
    \draw[arrow, dashed, gray] (x2.east) -- ++(0.5, 0) |- (add2.east);
\end{tikzpicture}
\caption{Post-normalization (left) vs Pre-normalization (right)}
\label{fig:pre-vs-post-norm}
\end{figure}

\subsection{Why Pre-Norm?}

\begin{enumerate}
    \item \textbf{Stable gradients}: The residual connection provides a direct gradient path without normalization in the way
    \item \textbf{Easier training}: Works with larger learning rates and simpler hyperparameter tuning
    \item \textbf{No warmup required}: Post-norm often needs learning rate warmup; pre-norm is more forgiving
\end{enumerate}

\subsection{Implementation}

\begin{lstlisting}[style=python, caption={TransformerBlock with pre-norm from model.py}]
class TransformerBlock(nn.Module):
    def __init__(self, config: ModelConfig):
        super().__init__()
        self.attn_norm = RMSNorm(config.d_model)
        self.ffn_norm = RMSNorm(config.d_model)
        self.attention = MultiHeadSelfAttention(config)
        self.ffn = FeedForward(config)

    def forward(self, x: Tensor, attention_mask: Tensor | None = None) -> Tensor:
        # Pre-norm attention with residual
        x = x + self.attention(self.attn_norm(x), attention_mask)

        # Pre-norm FFN with residual
        x = x + self.ffn(self.ffn_norm(x))

        return x
\end{lstlisting}

% -----------------------------------------------------------------------------
\section{Residual Connections}
% -----------------------------------------------------------------------------

\textbf{Residual connections} (skip connections) add the input directly to the output:
\begin{equation}
\text{output} = x + f(x)
\end{equation}

\subsection{Why Residual Connections?}

\begin{enumerate}
    \item \textbf{Gradient flow}: Gradients can flow directly through the addition, avoiding vanishing gradient problems
    \item \textbf{Identity as baseline}: The layer can learn to add corrections rather than complete transformations
    \item \textbf{Deep networks}: Enable training very deep networks (GPT-3 has 96 layers!)
\end{enumerate}

\subsection{Gradient Flow Visualization}

\begin{figure}[H]
\centering
\begin{tikzpicture}[
    box/.style={rectangle, draw, minimum width=2cm, minimum height=0.8cm},
    arrow/.style={-{Stealth[length=2mm]}, thick},
    grad/.style={-{Stealth[length=2mm]}, thick, red!70!black},
]
    % Without residual
    \node at (-3, 2) {\textbf{Without Residual:}};
    \node[box, fill=blue!10] (a1) at (-3, 0.8) {Layer 3};
    \node[box, fill=blue!10] (a2) at (-3, -0.6) {Layer 2};
    \node[box, fill=blue!10] (a3) at (-3, -2) {Layer 1};

    \draw[arrow] (a1) -- (a2);
    \draw[arrow] (a2) -- (a3);

    % Gradient
    \draw[grad] (-4, -2) -- (-4, 0.8);
    \node[font=\footnotesize, red!70!black, left] at (-4, -0.6) {Gradient};
    \node[font=\footnotesize, red!70!black, left] at (-4, 1.2) {may vanish};

    % With residual
    \node at (3, 2) {\textbf{With Residual:}};
    \node[box, fill=blue!10] (b1) at (3, 0.8) {Layer 3};
    \node[box, fill=blue!10] (b2) at (3, -0.6) {Layer 2};
    \node[box, fill=blue!10] (b3) at (3, -2) {Layer 1};

    \draw[arrow] (b1) -- (b2);
    \draw[arrow] (b2) -- (b3);

    % Skip connections
    \draw[arrow, dashed, gray] (3.8, 0.4) -- ++(0.5, 0) -- ++(0, -2.4) -- ++(-0.5, 0);

    % Gradient (direct path)
    \draw[grad, thick] (4.8, -2) -- (4.8, 0.8);
    \node[font=\footnotesize, red!70!black, right] at (4.8, -0.6) {Direct};
    \node[font=\footnotesize, red!70!black, right] at (4.8, 1.2) {gradient path};
\end{tikzpicture}
\caption{Residual connections enable direct gradient flow through deep networks}
\label{fig:residual-gradient}
\end{figure}

\begin{keypoint}
Without residual connections, gradients must pass through every layer. With residuals, gradients have a ``highway'' directly back to earlier layers, enabling much deeper networks.
\end{keypoint}

% -----------------------------------------------------------------------------
\section{Dropout in Transformer Blocks}
% -----------------------------------------------------------------------------

Dropout is applied in three places within a transformer block:

\begin{enumerate}
    \item \textbf{After attention weights}: Randomly drops attention connections
    \item \textbf{After FFN hidden layer}: Regularizes the feed-forward network
    \item \textbf{After attention output projection}: Before residual addition
\end{enumerate}

\begin{lstlisting}[style=python]
# In attention:
attn_weights = self.attn_dropout(attn_weights)  # After softmax

# In FFN:
hidden = self.dropout(hidden)  # After SwiGLU, before down_proj
\end{lstlisting}

\TryLM uses \texttt{dropout=0.1} by default.

% -----------------------------------------------------------------------------
\section{The Complete Transformer Block}
% -----------------------------------------------------------------------------

Let's trace data through a complete block:

\begin{lstlisting}[style=python]
# Input: [batch=2, seq_len=128, d_model=384]
x = ...

# TransformerBlock forward pass:

# 1. Pre-norm for attention
normed = self.attn_norm(x)  # [2, 128, 384]

# 2. Multi-head self-attention
attn_out = self.attention(normed, attention_mask)  # [2, 128, 384]

# 3. Residual connection
x = x + attn_out  # [2, 128, 384]

# 4. Pre-norm for FFN
normed = self.ffn_norm(x)  # [2, 128, 384]

# 5. Feed-forward with SwiGLU
ffn_out = self.ffn(normed)  # [2, 128, 384]

# 6. Residual connection
x = x + ffn_out  # [2, 128, 384]

# Output: [2, 128, 384] (same shape as input!)
\end{lstlisting}

\begin{keypoint}
Each transformer block preserves the input shape: \texttt{[batch, seq\_len, d\_model]}. This enables stacking arbitrary numbers of blocks.
\end{keypoint}

% -----------------------------------------------------------------------------
\section{Stacking Blocks}
% -----------------------------------------------------------------------------

A complete transformer stacks multiple blocks:

\begin{lstlisting}[style=python, caption={Stacking transformer blocks from model.py}]
class GPT(nn.Module):
    def __init__(self, config: ModelConfig) -> None:
        super().__init__()
        # ...

        # Stack of transformer blocks
        self.blocks = nn.ModuleList([
            TransformerBlock(config)
            for _ in range(config.n_layers)  # 6 blocks
        ])

    def forward(self, x: Tensor, ...) -> Tensor:
        x = self.token_embedding(input_ids)
        x = self.dropout(x)

        # Pass through all blocks
        for block in self.blocks:
            x = block(x, attention_mask)

        x = self.final_norm(x)
        # ...
\end{lstlisting}

For \TryLM with 6 layers:
\begin{itemize}
    \item Each block has $\sim$1.7M parameters
    \item 6 blocks $\times$ 1.7M = $\sim$10.2M parameters
    \item Plus embeddings ($\sim$6.3M) = $\sim$10.8M total
\end{itemize}

% -----------------------------------------------------------------------------
\section{Exercises}
% -----------------------------------------------------------------------------

\begin{exercise}
\textbf{Calculate RMSNorm}

Given vector $x = [2, 4, 6]$:
\begin{enumerate}
    \item Compute the mean square: $\frac{1}{3}(2^2 + 4^2 + 6^2)$
    \item Compute the RMS: $\sqrt{\text{mean square}}$
    \item Normalize: $x / \text{RMS}$
\end{enumerate}
\end{exercise}

\begin{exercise}
\textbf{Compare Activations}

Implement and compare different activations:
\begin{lstlisting}[style=python, numbers=none]
import torch
import torch.nn.functional as F
import matplotlib.pyplot as plt

x = torch.linspace(-3, 3, 100)

plt.plot(x, F.relu(x), label='ReLU')
plt.plot(x, F.gelu(x), label='GELU')
plt.plot(x, F.silu(x), label='Swish/SiLU')
plt.legend()
\end{lstlisting}

How do they differ around $x = 0$?
\end{exercise}

\begin{exercise}
\textbf{Count Block Parameters}

For a single transformer block with $d_{\text{model}} = 384$, $d_{\text{ff}} = 1536$, and 6 attention heads, count:
\begin{enumerate}
    \item Parameters in QKV projection
    \item Parameters in output projection
    \item Parameters in FFN (gate, up, down)
    \item Parameters in both RMSNorm layers
    \item Total parameters per block
\end{enumerate}
\end{exercise}

% -----------------------------------------------------------------------------
\section{Summary}
% -----------------------------------------------------------------------------

\begin{summary}
    \item A \textbf{transformer block} = attention + FFN + residuals + norms
    \item \textbf{RMSNorm} is simpler and faster than LayerNorm, with comparable performance
    \item \textbf{SwiGLU} uses gating: $\text{Swish}(\text{gate}) \odot \text{up}$, better than GELU
    \item \textbf{Pre-normalization} normalizes before sublayers, enabling stable training
    \item \textbf{Residual connections} enable gradient flow and very deep networks
    \item Blocks preserve shape: $[B, S, D] \to [B, S, D]$, allowing arbitrary stacking
\end{summary}

Now we have all the components of a transformer block. In the next chapter, we'll assemble the complete GPT model by combining embeddings, stacked blocks, and the output projection.
